% paper/framing_gamma.tex -- framing note: the labour axis as a gradient
% (BF <-> GAMMA) and the sectoral family that spans it.
%
% Companion documents: paper/equivalence.tex (v3, Section 5 carries the
% proposition and proof), docs/decisions/ADR-0022-sectoral-labour-markets.md
% (the decision), docs/WORKPLAN_SENSITIVE_PRICES.md (the executable annexe),
% paper/tables/matrix_5x3_v9_flows.md (the flows behind the numbers).
%
% Evidence: the matrix_5x3_v9 generation (33 cells, 15/15 matrix + 18 sectoral,
% all gates pass, 2026-09-19) and the matrix_5x3_v6 generation it reproduces
% bit-for-bit on the fifteen matrix cells.
%
% Version 1 (September 2026). No LaTeX in the analysis container: compile on
% the Mac (pdflatex; plain thebibliography below, no .bib dependency).

\documentclass[11pt]{article}

\usepackage[english]{babel}
\usepackage[a4paper,top=2.2cm,bottom=2.2cm,left=3cm,right=3cm]{geometry}
\usepackage{amsmath,amssymb,amsthm}
\usepackage{booktabs}
\usepackage{xcolor}
\providecolor{revisionV1}{HTML}{800080}
\providecolor{revisionV2}{HTML}{CC5500}
\providecolor{revisionV3}{HTML}{006400}
\providecolor{revisionV4}{HTML}{0000CD}
\providecolor{revisionV5}{HTML}{808080}
\providecolor{revisionV6}{HTML}{8B0000}
\usepackage[colorlinks=true,allcolors=blue]{hyperref}

\setlength\parindent{0pt}
\setlength{\parskip}{0.45em}

\newtheorem{proposition}{Proposition}
\newtheorem{remark}{Remark}

\newcommand{\Lbar}{\bar{L}}

\title{Framing the revision: the labour axis as a gradient\\[0.4em]
\large BF $\leftrightarrow$ GAMMA, and the sectoral family between them\\[0.2em]
\normalsize A framing note on the sectoral-labour-market result (ADR-0022)}
\author{Hermes Agent (Lt.\ Cmdr.\ Data), for Prof.\ Dr.\ J.\ Kapeller}
\date{September 19, 2026}

\begin{document}
\maketitle

\begin{center}\small
\textbf{Version 1} (September 2026)
\end{center}

\begin{abstract}
The executed $5\times3$ matrix carries two labour-rigidity rows at its
endpoints: BF, where the allocation is frozen and the $N$ sectoral wages absorb
the shock, and GAMMA, where the real wage is pinned and employment absorbs.
Between them the matrix is empty, and in every intermediate row prices are
pinned to the baseline to machine precision. The sectoral-labour-market closure
(ADR-0022) fills that gap: it makes the labour axis \emph{continuous}, with the
already-executed BF row as its rigid corner and a GAMMA-like regime as its
elastic limit. This note fixes what that licenses for the revision's framing,
what it obliges us to disclose, and what still has to be measured before a
number can be stated rather than a band. A further section supplies the missing
motivation for the one genuinely discretionary modelling choice --- which
sectors are rigid --- from capacity pressure in the sectors the programme
presses on, and reports the rigidity sweep as the robustness exhibit. A final
section tabulates the five GAMMA variants against their measured outcomes.
\end{abstract}

\section{The result, in one paragraph}

With one aggregate real-wage equation, demand reaches prices only through the
wage, and the price block is homogeneous of degree one in $(p,w)$ and
demand-free, so under a demand-only programme $p = w = \Pi = 1$ and the whole
quantity system collapses onto a linear Leontief multiplier
(\texttt{paper/equivalence.tex}, Lemma 1). That is why the matrix's
flexible-wage rows are uninformative about prices. Replacing the single
equation by $N$ sectoral labour markets,
\begin{equation}
L^{cm}_i(p, y_i, w_i) \;=\; \bar{L}_i\,
\Big[ \frac{w_i}{\Pi(p)} \Big]^{\eta_{s,i}},
\qquad i = 1, \dots, N, \quad \eta_{s,i} \geq 0,
\label{eq:family}
\end{equation}
ties each wage to that sector's own quantity, so the $N$ zero-profit conditions
acquire a demand dependence and $p = \pi(A, w)$ becomes demand-sensitive. Three
properties of the resulting family are exact, and all three are now proved and
measured:

\begin{itemize}
\item \textbf{The rigid corner \emph{is} the BF row.} With $\eta_{s,i} = 0$ for
      every sector the family reduces to the $\eta = 0$ endpoint of ADR-0020
      option C --- the same equations, not a limit of them
      (Proposition~3 of \texttt{paper/equivalence.tex}; measured
      $\max|\Delta| = 0.00\times10^{0}$ in all four blocks when warm-started
      from the endpoint solution, below $10^{-10}$ from a cold start). The
      mechanism is therefore anchored on an executed, gated, published set of
      cells, and no separate rigid sectoral cell is reported.
\item \textbf{The elastic end is GAMMA-like, but is not GAMMA.} As
      $\eta_{s,i} \to \infty$ the real wage returns to its technology-determined
      anchor and employment becomes the absorbing margin --- GAMMA's prices and
      GAMMA's absorbing margin. It is a limit, not a cell: the supply slope
      diverges and the direct solve fails ($\eta_s = 10^6$). It is also not the
      fixed-wage system itself, which pins $F \equiv 0$ while the family keeps
      the transfer free. The family \emph{spans} the two rigidity rows without
      terminating in either.
\item \textbf{The executed matrix is untouched.} The fifteen matrix cells were
      re-executed on the extended kernel (\texttt{matrix\_5x3\_v9}): every
      headline metric and diagnostic reproduces the v6 values bit-for-bit
      ($\max|\Delta| = 0$). The extension is provably non-invasive.
\end{itemize}

Table~\ref{tab:ladder} gives the executed evidence. The discriminator the
framing turns on is not ``do prices move'' but ``do they move \emph{with the
demand composition}'': in the sectoral cells $\max|p-1|$ differs between F1 and
F2/F3 at every rung, while all twelve $\eta = 1$ cells of the executed matrix
read exactly $0.000000$.

\begin{table}[htbp]
\centering
\caption{The sectoral family, executed (\texttt{matrix\_5x3-v9-BETA-*};
deviations from the no-programme baseline; F2 and F3 agree to all printed
decimals in every cell). The rigidity share matters more than the elasticity
level: the last two uniform rows and the two-group rows differ only in
\emph{which} sectors are rigid.}
\label{tab:ladder}
\begin{tabular}{lrrrr}
\toprule
Sectoral cell & $\max|p-1|$ & $\max|p-1|$ & Employment & Consumption \\
 & (F1) & (F2) & (F2) & (F2) \\
\midrule
uniform $\eta_s = 0.25$ & $0.133610$ & $0.152897$ & $1.000864$ & $-0.016888$ \\
uniform $\eta_s = 0.5$  & $0.095839$ & $0.105666$ & $1.001255$ & $-0.015712$ \\
uniform $\eta_s = 1$    & $0.061283$ & $0.065422$ & $1.001627$ & $-0.014669$ \\
uniform $\eta_s = 2$    & $0.035629$ & $0.037171$ & $1.001914$ & $-0.013912$ \\
\midrule
rigid programme sectors, $\eta_s = 0.5$ & $0.215061$ & $0.272033$ & $0.996621$ & $-0.026914$ \\
rigid largest half, $\eta_s = 0.5$ & $0.109583$ & $0.123362$ & $1.000196$ & $-0.019053$ \\
\midrule
BF ($\eta_{s,i} = 0$, ADR-0020) & $0.221442$ & $0.279009$ & $1.000000$ & $-0.019800$ \\
GAMMA (fixed wage, for contrast) & $0.000000$ & $0.000000$ & $1.001260$ & $-0.015333$ \\
\bottomrule
\end{tabular}
\end{table}

Two readings beyond the price column. First, the elasticity \emph{splits the
adjustment}: as $\eta_s$ rises the price response falls monotonically
($0.1529 \to 0.0372$ at F2) while employment rises ($1.000864 \to 1.001914$) ---
the family trades price adjustment for quantity adjustment, and the welfare
cost falls with it ($-1.69\% \to -1.39\%$). Second, and more consequentially,
the \emph{incidence} of the rigidity dominates its level: making the programme
sectors rigid more than doubles the price response ($0.2720$ against $0.1057$)
and \emph{reverses} the employment effect ($-0.34\%$ against $+0.13\%$), while
rigidity on the largest half instead leaves employment near neutral
($+0.02\%$) at $0.1234$.

\section{Five framing implications}

\begin{enumerate}
\item \textbf{The critique is answered constructively, not conceded.} The
      reviewer objection that the matrix's labour rows are a normalisation ---
      that BETA reproduces ALPHA and DELTA reproduces GAMMA, so the labour axis
      carries no information --- is correct \emph{about the single-wage
      parametrisation} and wrong about the mechanism. The sectoral family
      repairs it inside the BETA family, with no new closure id, no new
      taxonomy, and no new endpoint to justify: its rigid corner is the
      endpoint the paper has already accepted. The revision can therefore
      present the labour axis as informative and say precisely where the
      information comes from.
\item \textbf{The two-paradigm reading becomes a gradient.} The revision's
      current framing is a dichotomy: the flexible-wage class (composition at
      pinned employment, $F$ as the instrument) against the fixed-wage class
      (employment at pinned real wage, $F \equiv 0$). The sectoral family turns
      that dichotomy into a one-parameter path between the two classes: BF at
      the rigid corner (quantities frozen, prices free), GAMMA at the elastic
      end (prices frozen, quantities free), with the elasticity vector
      governing the split. This is the framing device of this note: the matrix's
      labour axis is not two boxes but a continuum of adjustment margins, and
      the financing closure remains the orthogonal axis. For a heterodox
      readership the gradient is the better story --- ``which margin absorbs''
      becomes a measurable continuum rather than a binary label --- and it is
      the same object the ETAs workplan calls the labour door.
\item \textbf{A nominal/real distinction exists in the $\eta = 1$ rows for the
      first time.} Every $\eta = 1$ cell of the executed matrix has a GDP
      deflator of exactly $1.000000$; the sectoral cells read $1.0010$ to
      $1.0063$. The revision can therefore discuss real versus nominal effects
      in the flexible-wage rows --- previously possible only in the fixed-wage
      rows and at BF --- which matters for the ``who pays'' argument: the
      sectoral rows show a real welfare cost that is not merely a tax
      reallocation.
\item \textbf{A disclosure obligation attaches to the numbers.} The labour
      specification moves the revision's headline welfare numbers by up to
      $2.7\%$ of consumption (rigid-programme variant) --- more than the
      GAMMA-versus-ALPHA difference on which the equivalence note is built ---
      and it moves them \emph{against} the employment direction in the rigid
      case (employment falls while the welfare index falls further). The
      specification is therefore not a technicality to be relegated to an
      appendix: whichever variant is reported, the sensitivity to the labour
      market's structure must be visible in the main text.
\item \textbf{An editorial choice remains open, and it is the operator's.}
      Whether the sectoral family enters the \emph{published} matrix as extra
      rows, or appears as the separate section that resolves the labour door
      (the form adopted in \texttt{paper/equivalence.tex} v3). The
      recommendation is the separate section: it keeps the $5\times3$ as the
      object the reviewer attacked, presents the family as the answer to it,
      and avoids committing the published table to a grouping rule that is
      still a modelling choice (item 2 of the next section).
\end{enumerate}

\section{What is still needed before a number can be stated}

Five items, ordered by what they change. The first three are the minimum for a
clear statement; the last two are polish.

\begin{enumerate}
\item \textbf{The level of $\eta_{s}$ is unidentified.} Everything above is a
      sensitivity band, not an estimate: under demand-only shocks the price
      response is monotone in $\eta_s$ and the data cannot choose a rung. Two
      routes close this. (a) The supply-shock arm: at $A \neq 1$ the real wage
      leaves its anchor and $\eta_s$ becomes identifiable from the
      wage--employment response (this is the same statement as the
      supply-side remark of \texttt{paper/equivalence.tex}, Section 4). (b) An
      external
      calibration: aggregate labour-supply elasticities in the literature sit
      around $0.1$--$0.5$, which brackets the middle of the ladder --- but that
      must be declared as an assumption, not presented as an estimate.
\item \textbf{The grouping rule is the weak link.} ``Rigid on the programme
      sectors'' and ``rigid on the largest half'' differ by a factor $2.5$ in
      the price response and \emph{flip the sign} of the employment effect.
      Either a principled rule is found (sectoral supply elasticities, hiring
      frictions, an institutional criterion) or the range across rules is
      reported as a band and the sign is presented as rule-dependent. This is
      the single most consequential open choice, because it determines the
      direction of the employment result. Measured since
      (\texttt{probe16\_rigidity\_share.jl}): the band is price response
      $[0.0969, 0.2790]$, employment $[-0.31\%, +0.16\%]$ and consumption
      $[-2.63\%, -1.50\%]$; and \emph{which} sectors are rigid dominates
      \emph{how many} --- the programme ranking saturates as soon as the seven
      programme sectors are inside the rigid group ($0.2731$ at a rigid share of
      $0.25$ against $0.2790$ at $1$), whereas rigid large employers move the
      price response only from $0.1057$ to $0.1262$. The recommendation is
      therefore concrete: report the band, and state the rule --- and
      Section~\ref{sec:capacity} supplies the missing motivation for the rule
      itself: capacity pressure in the sectors the programme presses on makes
      the rigid group a testable economic statement rather than a partition
      chosen to produce a result.
\item \textbf{The two arms must be combined.} All sectoral evidence is
      demand-only, while $\eta_s$ is only \emph{identified} under supply shocks.
      Running the sectoral closure under the supply-shock scenario decides
      whether the mechanism is a headline or a robustness remark: if the
      technology shock dominates the price response, the sectoral channel's
      contribution shrinks; if the two add up, the mechanism is a first-order
      part of the story. This is one probe and one design. Run since
      (\texttt{probe15\_sectoral\_supply.jl}), and the answer is mixed but
      informative: under a $+20\%$ shock in sector 1 the rungs \emph{separate}
      (employment $1.000000$, $1.002187$, $1.008218$ and the real-wage spread
      $1.215$, $1.105$, $1.043$ at $\eta_s = 0, 0.5, 2$), so the elasticity is
      identifiable there --- but the shock roughly doubles the price response
      and flips the sign of the welfare effect (supply-only $+1.20\%$, supply
      $+$ programme $-0.40\%$, demand-only $-1.57\%$), so the two channels do
      not add linearly and the technology shock dominates the price movement.
      One precision must travel with any supply-arm table: under a supply shock
      $\max|p-1|$ is not a demand-sensitivity measure (the shock alone moves
      prices by $0.226$ with no programme), so the demand channel is read from
      the F1/F2/F3 differences. A supply-shock \emph{design} (ADR-0018's schema
      item) is still owed.
\item \textbf{The welfare metric needs a caveat, and the caveat is structural.}
      The T\"ornqvist index on household consumption is dispersion-blind: it
      ignores the welfare cost of the relative-wage movement the mechanism
      generates ($\max_i w_i/\min_i w_i$ of $1.26$ at $\eta_s = 0.5$ and $1.70$
      in the rigid-programme variant, for a programme worth $1.33\%$ of GDP).
      The kernel cannot supply the missing part either: the labour supply has no
      leisure term and no income effect (DE-0004), so there is no apparatus in
      which the employment and relative-wage movement could be valued. The
      honest form is therefore the restrictive one --- report the consumption
      effect, and present employment and wage dispersion as \emph{allocation}
      facts rather than welfare claims.
\item \textbf{Polish.} A finer ladder ($\eta_s \in \{0.1, 5\}$) to fill the
      tails, and a robustness probe on the sign of the rigid-group employment
      effect under a different shock vector (a broader programme, or a uniform
      shock).
\end{enumerate}

\section{The grouping rule, motivated: capacity pressure in the shocked sectors}\label{sec:capacity}

The rigid group of Section~2's item (2) need not be an arbitrary partition. The
programme's incidence supplies a principled candidate, and the economics of
capacity supplies the reason to expect the candidate to be right.

\paragraph{The incidence is concentrated.} The 2024 impulse data place the whole
programme in \emph{seven} of the seventy-one sectors, and $63.4\%$ of it in a
single one: specialised construction works. The remaining six are material and
equipment sectors --- rubber and plastics ($9.4\%$), ceramic products and
processed stone ($8.0\%$), chemicals ($7.6\%$), glass ($6.1\%$), machinery
($4.3\%$), electrical equipment ($1.3\%$). A programme that is two-thirds one
construction sector and one-third its input suppliers is not a diffuse demand
stimulus; it is a demand push aimed at the economy's project-capacity chain.

\paragraph{Why that makes those sectors the rigid ones.} Capacity pressure is
the classical short-run constraint on demand-led expansion, and the
capacity-utilisation and Verdoorn traditions treat it as the margin where
demand-led growth meets its supply-side limit
\cite{verdoorn1949,kaldor1966,lavoie2014}. In the sectors that a programme
presses on --- construction trades and their material suppliers --- the
short-run margin is capacity, not price: skilled trades cannot be trained and
project capacity cannot be built within the horizon of the impulse, so an
increase in demand is met largely at the prevailing (or rising) price with
little employment response. That is precisely the sectoral labour market of
equation~(1) with $\eta_{s,i} = 0$: a vertical short-run supply curve. The
grouping rule is therefore not a partition chosen to produce a result; it is the
statement that \emph{the sectors the programme presses on are the sectors whose
labour cannot expand}, which is the capacity-pressure diagnosis in the model's
own language. Two consequences follow. First, the rule is falsifiable and
testable against sectoral evidence (vacancies, overtime, delivery times, order
backlogs in the construction trades) rather than assumed. Second, it links the
labour door to the capacity door of the companion note's scope section: the
rigid group is the set of sectors where capacity is the binding margin, which is
the same object the utilisation externality models on the price side.

\paragraph{The robustness exhibit.} Table~\ref{tab:band} reproduces the sweep
that measures how much the answer depends on the rule and on the rigid share
(\texttt{probe16\_rigidity\_share.jl}). It supports the motivated rule rather
than merely reporting a band.

\begin{table}[htbp]
\centering
\caption{The rigidity sweep. All cells are the sectoral family at
$\eta_s = 0.5$ for the flexible group, F2, deviations from the no-programme
baseline; the price column reports $\max|p-1|$ at F1 and F2. ``Programme
rank''$/$``employment rank'' order sectors by programme incidence$/$baseline
employment, descending; ``least exposed'' is the ascending programme order.}
\label{tab:band}
\begin{tabular}{lrrrr}
\toprule
Rigid group & $\max|p-1|$ (F1) & $\max|p-1|$ (F2) & Employment & Consumption \\
\midrule
none (uniform $\eta_s = 0.5$) & $0.095839$ & $0.105666$ & $1.001255$ & $-1.571\%$ \\
programme rank, share $0.25$ & $0.215939$ & $0.273076$ & $0.996911$ & $-2.626\%$ \\
programme rank, share $0.5$  & $0.216336$ & $0.273412$ & $0.997300$ & $-2.542\%$ \\
programme rank, share $0.75$ & $0.219801$ & $0.277006$ & $0.999370$ & $-2.084\%$ \\
programme rank, share $1.0$  & $0.221442$ & $0.279009$ & $1.000000$ & $-1.980\%$ \\
employment rank, share $0.5$ & $0.109583$ & $0.123362$ & $1.000196$ & $-1.905\%$ \\
least exposed, share $0.5$   & $0.096877$ & $0.106792$ & $1.001614$ & $-1.500\%$ \\
\bottomrule
\end{tabular}
\end{table}

Three readings, and they are what make the exhibit a robustness check rather
than a disclaimer.

\begin{enumerate}
\item \textbf{The price response saturates once the pressured sectors are
      inside the rigid group.} At a rigid share of only $0.25$ --- seven sectors
      of seventy-one --- the response is already $0.2731$ against $0.2790$ when
      the whole economy is rigid, and $0.1057$ when nothing is. The remaining
      three quarters of the economy add $0.006$ between them. So the choice that
      matters is whether the \emph{pressured} sectors are rigid, and the
      capacity argument answers exactly that question; the share is then almost
      irrelevant.
\item \textbf{The alternative rankings do not reach the same place.} Rigid
      large employers give $0.1234$ and rigid least-exposed sectors $0.1068$ at
      the same share --- an order of magnitude below the motivated rule. So the
      rule is doing real work, and the honest presentation is the one the
      capacity argument supports, with the sweep as the sensitivity: the price
      response is $[0.097, 0.279]$ across rules, and the motivated rule sits at
      the top.
\item \textbf{The employment sign is the fragile part, and the motivated rule
      also explains it.} Employment is non-monotone in the rigid share under the
      programme ranking ($1.001255$ / $0.996911$ / $0.997300$ / $0.999370$ /
      $1.000000$) and changes sign across rankings. Under the capacity reading
      this is not an embarrassment but a prediction: if the sectors the
      programme presses on cannot expand their labour, the programme's
      employment gain must come from elsewhere and can be offset by the
      contractionary effect of the induced price and wage movement. The
      employment result should therefore be reported as rule-dependent and
      small in absolute terms ($[-0.31\%, +0.16\%]$), while the price response
      --- the object the labour door was opened for --- is robust to the rule.
\end{enumerate}

\section{The GAMMA variants, and what each delivers}\label{sec:gammas}

The fixed-wage closure is not a single object. Seven variants are on the table,
and they differ in the one property the revision turns on: whether prices see
demand. Table~\ref{tab:gamma} reports the measured outcomes.

\begin{table}[htbp]
\centering
\caption{The GAMMA variants, measured (deviations from the no-programme
baseline). Executed baseline: \texttt{matrix\_5x3-v9-GAMMA-*}; pinned wage
structure: \texttt{probe14} (a $+10\%$ tilt on the programme sectors);
capacity channel and Kaldor--Verdoorn: \texttt{probe17} (the fixed-wage system
with $A^{eff}_i = (y_i/\lambda_i)^{-\delta}$ in the cost hook, solved by the
fixed-point iteration the reduced form implies; $\delta = 0$ reproduces the
executed cells exactly). ``Sees demand'' means $\max|p-1|$ differs across the
financing columns.}
\label{tab:gamma}
\begin{tabular}{llrrrrl}
\toprule
Variant & Fin. & $\max|p-1|$ & Deflator & Employment & Consumption & Sees demand \\
\midrule
baseline ($w \equiv 1$) & F1 & $0.000000$ & $1.000000$ & $0.999027$ & $-0.081\%$ & no \\
 & F2 & $0.000000$ & $1.000000$ & $1.001260$ & $-1.533\%$ & no \\
 & F3 & $0.000000$ & $1.000000$ & $1.017796$ & $+2.264\%$ & no \\
DELTA ($\equiv$ baseline) & all & $0.000000$ & $1.000000$ & $=$ baseline & $=$ baseline & no \\
pinned wage structure & all & $0.057820$ & --- & $1.006865$ & --- & no \\
\quad (identical in F1/F2/F3) & & & & & & \\
capacity channel, $\delta = +0.5$ & F1 & $0.098457$ & $1.003027$ & $1.004167$ & $+0.459\%$ & yes \\
 & F2 & $0.110761$ & $1.005379$ & $1.008489$ & $-1.082\%$ & yes \\
 & F3 & $0.128463$ & $1.021461$ & $1.038814$ & $+2.199\%$ & yes \\
Kaldor--Verdoorn, $\delta = -0.5$ & F1 & $0.123462$ & $1.003454$ & $0.997837$ & $-1.211\%$ & yes \\
 & F2 & $0.115629$ & $1.001283$ & $0.998527$ & $-2.464\%$ & yes \\
 & F3 & $0.128014$ & $0.985336$ & $1.003370$ & $+2.266\%$ & yes \\
\midrule
BF allocation rule, $\eta = 0$ & F1 & $0.000000$ & $1.000000$ & $1.000000$ & $+0.043\%$ & no \\
 & F2 & $0.000000$ & $1.000000$ & $1.000000$ & $-1.694\%$ & no \\
 & F3 & $0.000000$ & $1.000000$ & $1.000000$ & $+0.000\%$ & no \\
dual labour market, insiders rigid & F1 & $0.024701$ & $1.000438$ & $0.989569$ & $-1.338\%$ & yes \\
 & F2 & $0.027768$ & $1.000492$ & $0.991179$ & $-2.881\%$ & yes \\
 & F3 & $0.030189$ & $1.000648$ & $1.005996$ & $+0.689\%$ & yes \\
\bottomrule
\end{tabular}
\end{table}

\begin{enumerate}
\item \textbf{The baseline is price-free by construction, and that is its
      content.} The deflator is exactly $1.000000$ in every financing column, so
      the programme is absorbed entirely by quantities: employment
      ($-0.10\%$ under F1, $+0.13\%$ under F2, $+1.78\%$ under F3) and the
      welfare index ($-0.08\%$, $-1.53\%$, $+2.26\%$). No nominal/real
      distinction exists in this row. The economically interesting object is the
      \emph{financing} spread, which is wide: F3 is the Keynesian case and F2
      the tax-financed loss.
\item \textbf{DELTA is the same row,} and that is the definitional equivalence:
      in the fixed-wage gauge the CES elasticities $(\theta,\epsilon,\sigma)$
      are unidentified, because prices are pinned and the quantity system
      reduces to the Leontief multiplier.
\item \textbf{The pinned wage structure (ADR-0021) buys distribution, not
      demand.} Prices move by $0.058$ under a $+10\%$ tilt on the programme
      sectors, but \emph{identically} in F1, F2 and F3: the movement is carried
      by the exogenous pin, so the price block still never sees demand. What the
      variant does add is the allocation margin and an explicit distributional
      instrument --- relative wages of $1.10$ and employment $+0.69\%$ under the
      tilt. It is the right instrument for a question about who works where, and
      the wrong one for a question about price response.
\item \textbf{The capacity channel gives a demand-sensitive price response with
      the wage still pinned.} At $\delta = +0.5$ the price response is $0.0985$, $0.1108$
      and $0.1285$ across the financing columns --- a genuine demand channel
      with the wage still pinned, no labour-market assumption, and no wage
      dispersion. The sign of $\delta$ is the economics: $\delta > 0$ is
      capacity pressure (costs rise with demand, the F3 deflator reaching
      $1.0215$) and $\delta < 0$ is the Kaldor--Verdoorn case, where prices
      \emph{fall} with demand (the F3 deflator reads $0.9853$). Employment
      absorbs more than in the baseline (F3: $+3.88\%$ at $\delta > 0$), so the
      capacity channel shifts the burden from the price level to employment
      relative to the labour door. One caveat is structural: the price block
      stops being separable from quantities, so the system is a genuine
      simultaneous solve (the probe converges in $35$--$39$ iterations at
      $\delta = +0.5$ and in $70$--$76$ at $\delta = -0.5$, the slower rate
      being the positive feedback of increasing returns --- proximity to
      instability, not a numerical accident).
\item \textbf{The BF allocation rule in the fixed-wage row is the double-rigidity
      corner.} The BF closure's parameter \emph{is} the allocation elasticity
      $\eta$ (ADR-0010), and feeding it in at its rigid value gives the fixed
      wage \emph{and} the frozen allocation: with $w \equiv 1$ and
      $L_i = \bar{L}_i$, nothing can absorb the programme --- not the wage, not
      employment (a datum here, $L = 1.000000$ in every column), not the price
      level, since zero profit with $w = 1$ gives $p = 1$. The programme moves
      only the output composition ($\max_i|y_i/\lambda_i - 1| \approx 0.21$) and
      the tax mix; under F3 household consumption is \emph{literally unchanged}
      ($+0.000\%$), because the programme is externally financed and the
      household's real income is pinned on both sides. It is the benchmark the
      other variants are read against --- and the demonstration that the two
      rigidity rows of the matrix (BF and GAMMA) compose into a well-defined
      cell even though no cell of the matrix reports it. (The intermediate
      values $0 < \eta < 1$ are not available: they were retired with the B\&F
      reallocation wedge.)
\item \textbf{The dual labour market is the quantity-based counterpart of the
      sectoral family's price-based rigidity, and it behaves differently.} Here
      the wage stays pinned and the \emph{pressured sectors cannot expand their
      labour}: insiders sit at $\bar{L}_i$, outsiders at the cost-minimizing
      demand. The constraint binds through the \emph{cost} --- a sector that must
      produce $y_i$ with $\bar{L}_i$ can only substitute towards intermediates,
      and its constrained unit cost exceeds the unconstrained one --- so insider
      prices rise ($0.0278$ at F2) against $0.0015$ in the outsider sectors: a
      $19\times$ dualism signature, visible in the price vector rather than in
      the wages, which are identical in both segments. The bottleneck is then
      transmitted economy-wide through intermediate costs, and the effect is
      \emph{contractionary}: total employment falls under F1 and F2
      ($-1.04\%$, $-0.88\%$) and consumption falls ($-1.34\%$, $-2.88\%$),
      whereas external financing (F3) expands both ($+0.60\%$, $+0.69\%$). This
      is the dualism result proper: an inflexible sector that is pressed on
      raises everyone's costs, so the programme's employment gain must come from
      outside the pressured chain and can be more than offset. Two further
      readings. First, the constrained sectors sit at about half of their
      capacity ceiling ($y^{max}_i = A_i \alpha_i^{1/(\epsilon-1)}\bar{L}_i$,
      finite because labour and intermediates are complements at
      $\epsilon = 0.5$; the baseline sits at the factor share $\alpha_i$ of it),
      so the constraint bites through the substitution margin rather than at a
      hard ceiling. Second, the comparison with the price-based rigidity is
      instructive: the sectoral family's rigid-programme variant reads
      $-0.34\%$ employment and $-2.69\%$ consumption at F2, i.e. \emph{milder}
      on employment than the quantity-constrained dual ($-0.88\%$), because a
      vertical supply curve still lets the wage --- and through it the price ---
      absorb, whereas a quantity constraint forces the cost up. The two
      rigidities are therefore not interchangeable: one prices the bottleneck,
      the other rations it.
\item \textbf{The sectoral family's elastic limit is not a variant at all.} As
      $\eta_s$ grows the price response falls toward zero ($0.0084$ at
      $\eta_s = 10$) and the direct solve fails at $10^6$: it is the same
      outcome as the baseline in prices, approached but never attained, and it
      keeps the transfer $F$ free where the fixed-wage system pins $F \equiv 0$.
      It belongs in the framing as the elastic end of the gradient, not in a
      menu of choices.
\end{enumerate}

The choice, then, is not between seven flavours of the same thing. If the
question is \emph{which margin absorbs the programme}, the baseline and DELTA
answer it (employment), and the sectoral family answers it with prices. If the
question is \emph{who works where}, ADR-0021's pin is the instrument. If the
question is \emph{whether demand moves prices inside the fixed-wage paradigm},
the capacity channel and the dual labour market both do it --- the first through
cost, the second through rationing, and only the second is contractionary. The
capacity channel does it with the wage still pinned, which is the cleanest
separation of the price channel from the wage channel the model admits. The two additions sharpen the menu in opposite directions: the BF
allocation rule is the corner where nothing adjusts, and the dual labour market
is the variant where the *quantity* margin is rationed --- the one fixed-wage
setting in which the programme reduces employment.

\section{What the framing does not license}

\begin{itemize}
\item \textbf{$\eta_{s,i}$ is not an estimated labour-supply elasticity.} It is
      a sensitivity parameter of a reduced-form supply curve, and under
      demand-only shocks it is unidentified. No statement about actual labour
      supply behaviour may be attached to the ladder.
\item \textbf{A sectoral elasticity difference is not a statement about
      preferences or leisure.} Making a sector rigid is a modelling assumption
      about its labour market, not evidence that its workers are less
      responsive.
\item \textbf{The wage dispersion is an outcome, not a policy instrument.} The
      model has no wage policy; the dispersion is what clearing $N$ sectoral
      markets implies at the programme's incidence.
\item \textbf{``Demand-sensitive prices'' does not mean every regime.} The
      scalar rows remain demand-free \emph{by design} --- that is the content
      of the companion note's equivalences --- and the framing must not blur the
      two claims.
\end{itemize}

\section*{Revision log}

\noindent\textbf{Version 1} (September 2026) --- Initial framing note. Fixes
the five implications of the sectoral-labour-market result (ADR-0022) for the
revision's framing, the five items still needed before a number can be stated,
and the four claims the framing does not license. A section motivates the
rigid-group choice from capacity pressure in the sectors the programme presses
on (seven sectors, $63.4\%$ specialised construction) and reproduces the
rigidity sweep (rigid share and ranking) as the robustness exhibit. A further
section sets out the five GAMMA variants with their measured outcomes (the
capacity channel measured for the first time by \texttt{probe17}: the only
fixed-wage variant whose prices see demand). Evidence:
\texttt{matrix\_5x3\_v9} (33 cells), the bit-identity with
\texttt{matrix\_5x3\_v6} on the fifteen matrix cells,
\texttt{probe15\_sectoral\_supply.jl} and \texttt{probe16\_rigidity\_share.jl}.

\begin{thebibliography}{9}

\bibitem{bh2019}
Baqaee, D.\,R., and E.\ Farhi (2019). Beyond Hulten's theorem.
\emph{Econometrica} 87(3).

\bibitem{bf2019}
Baqaee, D.\,R., and E.\ Farhi (2020). Productivity and misallocation in
general equilibrium. \emph{Econometrica} 88(1).

\bibitem{robinson2006}
Robinson, S.\ (2006). Macro models and multipliers: Leontief, Stone, Keynes,
and CGE models. IFPRI, Washington, DC.

\bibitem{verdoorn1949}
Verdoorn, P.\,J.\ (1949). Fattori che regolano lo sviluppo della produttivit\`a
del lavoro. \emph{L'Industria} 1, 3--10.

\bibitem{kaldor1966}
Kaldor, N.\ (1966). \emph{Causes of the Slow Rate of Growth of the United
Kingdom}. Cambridge University Press, Cambridge.

\bibitem{lavoie2014}
Lavoie, M.\ (2014). \emph{Post-Keynesian Economics: New Foundations}.
Edward Elgar, Cheltenham.

\end{thebibliography}

\end{document}
